\setcounter{footnote}{20}

\printedpage{49}
% head: A. Grothendieck : Dérivateurs, Ch. VII, Édition des Universités de Montpellier II et Paris VII
\msp{page 74}
of the invertible arrows of $X'$ (which set does indeed satisfy the two admissibility conditions). Thus by what precedes we have a commutative diagram
\[
\begin{tikzcd}
X \ar[r,"\kappa"] \ar[d,"f"'] & \Pi_{\Sigma}(X) \ar[d,"\bar f = \Pi_{\Sigma}(f)"] \ar[dl,dotted,"\tilde f"'] \\
X' \ar[r,"\kappa'","\sim"'] & \Pi_{\Sigma'}(X'),
\end{tikzcd}
\]
where $\bar f$ is uniquely determined. Since $\kappa'$ is an isomorphism (as we saw earlier, since $\Sigma'$ \texttt{[is]} made up of the invertible arrows), the existence and uniqueness of $\bar f$ is equivalent to the announced conclusion,

\hfill q.e.d.

\textbf{One more sorites} on the homotopy classes of elements of a category $X$. In general, one sees that $X$ and $\Pi_{\Sigma}(X) = X\Sigma^{-1}$ have the same connected components, for if
\msp{page 75}
$a$ and $b$ in $X$ are joined by an arrow of $X\Sigma^{-1}$, i.e. by a $\Sigma$-path, then they are in the same connected component of $X$. But the converse is also true, as one sees e.g. by taking $\Sigma = \mathrm{Fl}(X)$, and by saying that in the \textit{groupoid} $\Pi(X)$, where two vertices are in the same \texttt{[connected]} component if and only if $\exists\ x \longrightarrow y$ (a relation which is not only transitive and reflexive, but also symmetric in a groupoid). A path
\[
a \xrightarrow{\quad c \quad} b
\]
will also be called a \textit{homotopy} from $a$ to $b$. We say that $a$ and $b$ are \textit{homotopic} if there exists a homotopy $c$ from $a$ to $b$, in which case there also exists an ``inverse'' homotopy $c^{o}$ from $b$ to $a$. This also means that $a$ and $b$ are in the same connected component of $X$, i.e. define the same element of $\pi_0(X)$.

The homotopy will be called \textit{simple} if $c$ is of type $\Delta_{+}$ or $\Delta_{-}$, i.e. if it is given by an arrow $u : a \longrightarrow b$ or $v : b \longrightarrow a$. If one is looking for a homotopy between $a$ and $b$, one will try to find it as ``short'' as possible. If one has a homotopy $c$ of length $l$ from $a$ to $b$, then on replacing $\tau$ by $\tau_{\red}$, hence $I_{\tau}$ by $I_{\tau\,\red}$, the corresponding reduced
\msp{page 76}
subinterval, one finds a homotopy from $a$ to $b$ of type $\tau_{\red}$, whose length is equal to the \textit{``useful length''} of the type $\tau$ of $c$ (which justifies the terminology). Let us note that if
\[
\tau = \Delta_{\varepsilon}^{l_1}\Delta_{-\varepsilon}^{l_2}\Delta_{\varepsilon}^{l_3}\cdots\Delta_{(-1)^{r}\varepsilon}^{l_r},
\]
then
\[
\tau_{\red} = \Delta_{\varepsilon}\Delta_{-\varepsilon}\Delta_{\varepsilon}\cdots\Delta_{(-1)^{r}\varepsilon},
\]
% foot: Version 22 décembre 2001
\printedpage{50}
% head: A. Grothendieck : Dérivateurs, Ch. VII, Édition des Universités de Montpellier II et Paris VII
\texttt{[rather twice} $(-1)^{r-1}$\texttt{]} where $r$ is the useful length of $\tau$ (the not useful length being $n = \sum\limits_{i\in[1,r]} l_i$).

When one starts with a homotopy between \textit{functors}
\[
X \longrightarrow Y,
\]
it is understood that we are dealing with homotopies in the category $\uHom(X,Y)$. (Or, depending on the context, in a specified subcategory of $\uHom(X,Y)$.)

\msp{page 77}
On page 61, I invoked the existence of a \textit{functor} of \textit{``composition of paths''}
\[
\Ch_{\Sigma}(X;a,b)\times\Ch_{\Sigma}(X;b,c) \ \longrightarrow\ \Ch_{\Sigma}(X;a,c),
\]
which on the fiber of the first member over $(\tau,\tau') \in \operatorname{Ob}(\Ty\times\Ty)$, would be given by
\[
\Ch_{\Sigma,\tau}(X;a,b)\times\Ch_{\Sigma,\tau'}(X;b,c) \ \longrightarrow\ \Ch_{\Sigma,\tau\tau'}(X;a,c), \tag*{$(*)$}
\]
the composition of paths of specified types $\tau$, $\tau'$. That we have here a \textit{functor,} not just a map on the sets of objects, is fairly obvious, and can be expressed for example with the help of the amalga\texttt{[ma]}tion diagram
\[
\begin{tikzcd}[column sep=small]
 & \varepsilon = \{0,1\} \ar[dl,hook'] \ar[dr,hook] & \\
I_{\tau} \ar[dr,hook] & & I_{\tau'} \ar[dl,hook'] \\
 & I_{\tau\tau'} &
\end{tikzcd}
\]
by transforming it by $\uHom(\cdot,X)$ (instead of transforming it by $\mathrm{Hom}_{\mathrm{Cat}}(\cdot,X)$ only), at least in the particular case where $\Sigma = \mathrm{Fl}(X)$ (which was the one on page 61). But the passage to arrows is fairly obvious - for two composites $c'_1\circ c_1$ and $c'_2\circ c_2$, from
\[
\begin{tikzcd}[column sep=small]
a \ar[rr,"c_1"] \ar[d,equal] & {} \ar[d,"w"'] & b \ar[rr,"c'_1"] \ar[d,equal] & {} \ar[d,"w'"] & c \ar[d,equal] \\
a \ar[rr,"c_2"] & {} & b \ar[rr,"c'_2"] & {} & c
\end{tikzcd}
\]
\texttt{[one]} sees that the homomorphisms $w''$ from the first to the second, in $\Ch_{\Sigma}(X;a,c)$, which are the identity on the vertex of index $n+1$ of the two paths (where $n = \operatorname{long}\tau$), corresponding to $1_{\tau}$ and to $0_{\tau'}$, ``are'' exactly the pairs
% foot: Version 22 décembre 2001
\printedpage{51}
% head: A. Grothendieck : Dérivateurs, Ch. VII, Édition des Universités de Montpellier II et Paris VII
\msp{page 78}
formed of a homomorphism $w : c_1 \longrightarrow c_2$ and $w' : c'_1 \longrightarrow c'_2$ in $\Ch_{\Sigma,\tau}(a,b)$ and $\Ch_{\Sigma,\tau'}(b,c)$ respectively, and if $w''$, $w$, $w'$ correspond in this way, we shall set
\[
w'' = w * w'.
\]
There is an obvious transitivity to check to see that this defines a \textit{functor} $(*)$. We can express the situation with the help of the diagram
\[
\begin{tikzcd}[column sep=small]
 & \Ty\times\Ty \ar[rrr,"\pi\ (\text{composition})"] \ar[dl,"p_1"'] \ar[dr,"p_2"] & & & \Ty \\
\Ty & & \Ty, & &
\end{tikzcd}
\]
where $p_1$, $p_2$ are the two projections. We consider over $\Ty\times\Ty$ the three fibered categories % [sic] fibrés -> fibrées, agreement with "catégories"
\[
p_1^*(\CH_{\Sigma,\tau}(a,b)),\ \ p_2^*(\CH_{\Sigma,\tau'}(b,c)),\ \ \pi^*(\Ch_{\Sigma,\tau\tau'}(a,c))
\]
\texttt{[rather} $\pi^*(\CH_{\Sigma,\tau\tau'}(a,c)$\texttt{]}, whose fibers at $(\tau,\tau')$ are respectively % [sic] \pi^*(\CH_{\Sigma,\tau\tau'}(a,c) -> \pi^*(\CH_{\Sigma,\tau\tau'}(a,c)), one closing parenthesis missing
\[
\Ch_{\Sigma,\tau}(a,b),\ \ \Ch_{\Sigma,\tau'}(b,c),\ \ \Ch_{\Sigma,\tau\tau'}(a,c).
\]
(We write $\CH$ instead of $\Ch$, to indicate that we regard the categories under consideration as (fibered) categories \textit{over} $(\Ty)$ (\footnote{But for $p_1^*$, $p_2^*$ and $\pi^*$ these are ordinary Cartesian products in Cat, there is no need to be afraid and to write $\CH$'s \ldots}).) Over $\Ty\times\Ty$, we have a (Cartesian) homomorphism of fibered categories
\[
p_1^*(\CH_{\tau}(a,b))\times_{\Ty\times\Ty} p_2^*(\CH_{\tau}(b,c)) \ \longrightarrow\ \pi^*(\CH_{\tau}(a,c)) \tag*{$(**)$}
\]
\texttt{[rather} $p_1^*(\CH_{\tau}(a,b))\times_{\Ty\times\Ty} p_2^*(\CH_{\tau'}(b,c)) \ \longrightarrow\ \pi^*(\CH_{\tau\tau'}(a,c))$\texttt{]},
\msp{page 79}
for it is OK on the fibers ($(*)$ page 77), and if
\[
(\varphi,\varphi') : (\tau_1,\tau'_1) \longrightarrow (\tau,\tau') \qquad (\varphi : \tau_1 \longrightarrow \tau,\ \varphi' : \tau'_1 \longrightarrow \tau')
\]
is an arrow in $\Ty\times\Ty$, we deduce from it a commutative diagram
\[
\begin{tikzcd}
\Ch_{\Sigma,\tau}(a,b)\times\Ch_{\Sigma,\tau'}(b,c) \ar[r,"\text{comp.}"] \ar[d,"\varphi^*_{a,b}\times(\varphi'_{b,c})^*"] & \Ch_{\Sigma,\tau}(a,c) \ar[d,"(\varphi*\varphi')^*_{a,c}"] \\
\Ch_{\Sigma,\tau_1}(a,b)\times\Ch_{\Sigma,\tau'_1}(b,c) \ar[r,"\text{comp.}"] & \Ch_{\Sigma,\tau_1\tau'_1}(a,c)
\end{tikzcd}
\]
% foot: Version 22 décembre 2001
\printedpage{52}
% head: A. Grothendieck : Dérivateurs, Ch. VII, Édition des Universités de Montpellier II et Paris VII
\texttt{[rather} $\Ch_{\Sigma,\tau\tau'}(a,c)$\texttt{]}, i.e. the composition of paths with fixed endpoints is compatible with the inverse images of paths under arrows between the types of paths.

But the Cartesian homomorphism $(**)$ of fibered categories over $\Ty\times\Ty$ is not quite the same thing as an innocuous functor
\[
\Ch_{\Sigma}(a,b)\times\Ch_{\Sigma}(b,c) \ \longrightarrow\ \Ch_{\Sigma}(a,c). \tag*{$(*)$}
\]
The formula $(*)$ defines in any case a map on the sets of objects, since the sets of objects of each of the three categories under consideration decompose as a sum of sets $\Ch_{\Sigma,\tau}(\bullet,\bullet)$, where $\tau$ runs through $\mathrm{Ty} = \operatorname{Ob}\Ty$. But does this extend in a natural way to morphisms?

\msp{page 80}
But now I remember that this famous Cartesian product $p_1^*(\text{\textit{fourbi} (gear)})\times p_2^*(\text{thingamajig})$ is none other than the plain old Cartesian product
\[
\Ch_{\Sigma}(a,b)\times\Ch_{\Sigma}(b,c),
\]
sent to $\Ty\times\Ty$ by the Cartesian product of the two projection functors
\[
\begin{aligned}
p_{a,b} : \Ch_{\Sigma}(a,b) &\ \longrightarrow\ \ \Ty\\
p_{b,c} : \Ch_{\Sigma}(b,c) &\ \longrightarrow\ \ \Ty,
\end{aligned}
\]
hence
\[
\begin{tikzcd}
\Ch_{\Sigma}(a,b)\times\Ch_{\Sigma}(b,c) \ar[d,"p_{a,b}\times p_{b,c}"] \\
\Ty\times\Ty
\end{tikzcd}
\]
\msp{page 81}
which is well and truly fibrant (as a Cartesian product of two fibrant functors) and has the right fibers! So we do find, by the general sorites on fibrant functors (here over $\Ty\times\Ty$) a functor of fibered categories over $\Ty\times\Ty$
\[
\Ch_{\Sigma}(a,b)\times\Ch_{\Sigma}(b,c) \ \longrightarrow\ \qquad \underbrace{\pi^*(\CH_{\Sigma}(a,c))}_{=\,(\Ty\times\Ty)\times_{\Ty}\Ch_{\Sigma}(a,c)\ \ \text{(\footnotemark)}} \qquad .
\]
\footnotetext{fiber product along $\pi : \Ty\times\Ty \longrightarrow \Ty$.}
Since on the other hand we have a functor from the last member (fiber product) to the factor $\Ch_{\Sigma}(a,c)$, we win. (We would have found it again by the direct argument I was about to make, and which I crossed out - but I am happier this way!)

% foot: Version 22 décembre 2001
\printedpage{53}
% head: A. Grothendieck : Dérivateurs, Ch. VII, Édition des Universités de Montpellier II et Paris VII
To do things properly, we would have to prove the \textit{associativity} of the composition functors, not just at the level of objects, by arguments like this, but it comes down simply to a relation of the type
\[
(w''*w')*w = w''*(w'*w)
\]
for a system of three arrows in
\[
\Ch_{\Sigma}(a,b),\ \Ch_{\Sigma}(b,c)\ \ \text{and}\ \ \Ch_{\Sigma}(c,d)
\]
respectively. It can only be trivially
\msp{page 82}
true, but tedious to spell out.

The cold sweats are over!

In summary, the formalism of the $\Ch_{\Sigma}(X)$ and of the $\Ch_{\Sigma}(X;a,b)$ developed up to now gives complete satisfaction on the following points:

\begin{enumerate}
\item[a)] Good formalism for the composition of paths (which was disgusting before).

\item[b)] Elegant treatment of the construction of categories of fractions $\mathcal{M}\Sigma^{-1}$. But I expect that it will be less pretty than before, when it comes to replacing an $(\mathcal{M},\Sigma)$ by $(\ \underbrace{\mathcal{M}(X)}_{=\,\uHom(X,\mathcal{M})}\ ,\Sigma_X)$'s ($X$ an index category) \texttt{[rather} $\uHom(X^{o},\mathcal{M})$\texttt{]}. It is to remedy this that I will have to develop comparison statements, between the category $\Ch_{\Sigma}$ constructed here, and others whose type in Hot lends itself to computations in a more direct way, as in the original state of the question. Moreover, we will surely have to turn to Ind-categories, in order to be able to follow the variance of $\mathcal{M}(X)\Sigma_X^{-1}$ for variable $X$.

\item[c)] Good formalism for ``homotopies'' and for compositions of homotopies (which is another way of speaking of paths and of their composition).
\end{enumerate}
% foot: Version 22 décembre 2001